
Việc xác định các thông số cơ bản của Trái Đất như bán kính và khối lượng là một trong những thành tựu quan trọng của vật lý cổ điển. Điều đáng chú ý là những đại lượng này có thể được xác định thông qua các thí nghiệm và quan sát tương đối đơn giản.

Trong bài toán dưới đây, ta sẽ khảo sát hai phương pháp tiêu biểu: phương pháp thiên văn học dựa trên hiện tượng Mặt Trời lặn và thí nghiệm Cavendish sử dụng cân xoắn để xác định hằng số hấp dẫn.
\subsection*{Đo bán kính Trái Đất}
Tại một bãi biển gần xích đạo, một học sinh muốn ước lượng bán kính \(R\) của Trái Đất. Biết rằng Trái Đất quay đều và hoàn thành một vòng quanh trục trong thời gian \(T = 24\) giờ.  

Học sinh quan sát Mặt Trời ở hai độ cao so với mực nước biển là \(h_1\) và \(h_2\) với \(h_2>h_1\). Khi Mặt Trời vừa lặn ở độ cao \(h_1\), đồng hồ được bắt đầu bấm. Học sinh di chuyển nhanh lên độ cao \(h_2\) và vẫn quan sát thấy Mặt Trời chưa lặn. Khi Mặt Trời lặn lần tiếp theo ở độ cao \(h_2\), đồng hồ dừng bấm. Thời gian đo được là \(\Delta t\).  

Giả sử khúc xạ khí quyển là không đáng kể và tia sáng từ Mặt Trời được coi là chùm song song.

\begin{enumerate}
\item Chứng minh rằng bán kính Trái Đất có thể tính theo công thức:
\[
R = \frac{2 \, (\sqrt{h_2} - \sqrt{h_1})^2}{\left( 2 \pi \frac{\Delta t}{T} \right)^2}.
\]

\emph{Áp dụng số:}
\[
h_1 = 2\,\text{m}, h_2 = 18\,\text{m}, \Delta t = 21{,}79\,\text{s}, T = 24\,\text{h}
\]

\item \emph{Gợi ý:} Đối với góc nhỏ \(\theta\), có thể xấp xỉ \(\tan\theta \approx \theta\) (tính bằng radian).
\end{enumerate}

\vspace{8pt}

\subsection*{Đo khối lượng Trái Đất (thí nghiệm Cavendish)}

\textbf{Dụng cụ thí nghiệm:}
\begin{itemize}
    \item Cân xoắn
    \item Thước chia độ
    \item Một thanh nhẹ dài \(2L\)
    \item Hai quả cầu nhỏ có khối lượng \(m\) và hai quả cầu lớn có khối lượng \(M\)
    \item Gương phẳng và một nguồn sáng điểm
    \item Màn quan sát
    \item Các sợi dây mảnh chịu lực tốt
\end{itemize}

\begin{enumerate}
\setcounter{enumi}{2}
\item Xác định phương pháp đo góc lệch nhỏ \(\theta\) từ các dụng cụ trên. Nêu ưu điểm và nhược điểm của phương pháp.

\item Trình bày cơ sở lý thuyết và các bước thực hành để xác định hằng số hấp dẫn \(G\) và khối lượng Trái Đất \(M_E\). Lực tương tác hấp dẫn giữa hai vật khối lượng \(m_1\) và \(m_2\) cách nhau khoảng \(r\) được xác định bởi:
\[
F = G \frac{m_1 m_2}{r^2},
\]
trong đó \(G\) là hằng số hấp dẫn.  

\emph{Gợi ý:} Sử dụng công thức cân bằng momen lực. Momen đàn hồi của dây xoắn được cho bởi:
\[
M = k \theta,
\]
trong đó \(k\) là hằng số đàn hồi của dây, \(\theta\) là góc lệch của thanh cân xoắn so với vị trí ban đầu.
\emph{Áp dụng số:}
\end{enumerate}

\begin{flushright}
    (biên soạn bởi lax và KhorzX)
\end{flushright}